\section*{Appendix Q. Trainable perturbation-identification pilots}
\label{sec:trainable-cell}
\setcounter{table}{0}
\renewcommand{\thetable}{Q.\arabic{table}}
\renewcommand{\theHtable}{Q.\arabic{table}}

\paragraph{Task and scope.}
Given an observed transcriptional response and five candidate interventions, the predictor selects the intervention that generated the response. We implement three strata: single-gene perturbations (VCC H1-hESC), double-gene overexpression (Norman), and chemical perturbations (Tahoe). These replace frozen-policy replay as the forward development target, but do not retroactively change any earlier scorer or result. Every stratum has ten disjoint simulated clients. All arms within a stratum use identical held-out queries and candidate rosters, sampled from identities before training. The primary metric is intervention-macro Top-1; micro Top-1, MRR, Top-3 and cosine-logit cross-entropy are secondary. The latter uses a fixed temperature of 0.1. The analytical uniform-random Top-1 expectation is 20\%.

\paragraph{Data and grouping.}
VCC uses 60/20/20 training/development/test interventions: the first two are a fixed split of the original train80, while the original validation20 supplies the retrospective pilot test. Each client owns six training and two development interventions. Clients are balanced by intervention count, not raw cell count (5,204--12,158 training cells per client). Training targets average all available cells within an intervention. Evaluation averages disjoint blocks of at most 64 cells, dropping residual blocks smaller than 16; 360 test queries represent 20 interventions, not 360 independent perturbations. Common untreated controls are explicitly a public reference. VCC targets are differences of mean log1p counts normalized to 10,000.

The official Norman release supplies 117 double-gene interventions, split 70/20/27 by whole gene pair; clients own seven training and two development pairs. We average the released profiles within each training pair. Test queries are 1,350 released scVI-derived differential-expression profiles for 27 pairs, not independently measured single cells. The task retains the release's log-fold-change target scale. The common untreated reference is constructed from designated public training-pool controls; perturbed responses remain private during fitting.

The Tahoe pilot uses 307 unique canonical chemical structures at context CVCL\_0023 and 5\,$\mu$M. The context/dose combination is chosen by compound coverage from metadata, not prediction scores; alias structures are deduplicated before splitting. Whole-drug roles contain 180/60/67 compounds, with 18 training and six development drugs per client. All source conditions are from the previously designated training pool, excluding legacy held-out targets. Each query is a condition-level mean response, and all five options share context and dose. The 307 source conditions aggregate 789,787 cells; this is not the number of training examples.

\paragraph{Public initialization and trainable interface.}
We load the public CELLxGENE-pretrained scDEBART checkpoint, not target-finetuned UMM weights or saved test predictions \citep{sung2026scdebart}. Its SHA256 is verified before strict state-dictionary loading. Official architecture definitions are pinned to commit \texttt{8996c7ca}; attention behavior is preserved. Fresh candidate-conditioned hidden states are computed from public controls and intervention identities. For genetic tasks, the input panel comprises 1,024 genes selected using public-control variation plus all intervention genes; every directly perturbed gene is removed from both loss and ranking coordinates. Inhibition/overexpression inputs use the released $-10/+10$ perturbation encoding. The backbone is frozen; its actual response MLP is trained, optionally with a zero-output-initialized width-32 residual adapter. This is head adaptation, not full-backbone fine-tuning or the complete original scDEBART fine-tuning recipe.

The released genetic interface does not condition on chemicals. We therefore add an explicit experimental chemical-conditioned head for Tahoe. A seeded, fixed 128-dimensional projection of unit-normalized Morgan-2048 fingerprints is concatenated with frozen untreated scDEBART gene features. The head copies the public parameters and initializes added input columns to zero; these columns subsequently train with the head. A panel of 512 mapped, testable genes is selected using public control expression only. This adapter is not claimed as native scDEBART drug support. Before adaptation, its predictions are identical across drugs; its realized tie-broken baseline therefore fluctuates around chance.

\paragraph{Local training and research.}
Ten persistent workers open only their own training and development response shards. Each FedAvg round broadcasts the shared head, runs one complete local epoch with batch size two interventions, and averages returned parameters weighted by local training-intervention count. AdamW resets each round, gradients are clipped at norm one, and each candidate starts from the same public initialization with seed 42. Twenty rounds define this first-pass schedule; convergence has not been established. Untreated reference features are public, while offline centralized data packaging simulates private institutions. Training uses process-local data access, not an adversarial filesystem sandbox, secure aggregation or differential privacy.

The fixed recipe uses learning rate $10^{-4}$, weight decay $10^{-3}$ and MSE. A constrained research prototype permits learning rates $\{10^{-5},3\times10^{-5},10^{-4},3\times10^{-4}\}$, weight decays $\{0,10^{-3},10^{-2}\}$, MSE/weighted-MSE/MSE-plus-cosine objectives, and the optional residual adapter. Local Qwen2.5-7B-Instruct proposes a hypothesis, configuration, experiment and expected effect. Direct and loop arms share this checkpoint, prompt contract, three proposal slots and training schedule. Direct proposals receive no previous score feedback; loop proposals receive the complete aggregate development history. Both select by macro Top-1, then MRR, then lower cross-entropy. Thus this prototype isolates feedback access under bounded head search; it is not a rerun of the earlier native source-editing controller. Repeated configurations still consume their slots. Client 0 is the preselected one-laboratory control, not the weakest client chosen after scoring. All selections are sealed before test evaluation.

\begin{table}[htbp]
\centering\small
\caption{First-pass five-option perturbation identification. Entries are intervention-macro Top-1 (\%). Bold marks the highest measured value in each column, including ties; it does not imply SOTA or significance. Uniform random is an analytical expectation.}
\label{tab:trainable-cell-pilot}
\begin{tabular}{lrrr}
\toprule
Method & VCC single & Norman double & Tahoe drug \\
\midrule
Public head + cosine & \textbf{14.03} & 17.85 & 20.90 \\
1 lab, fixed head & 13.14 & 18.52 & 13.43 \\
1 lab, Qwen direct & 13.14 & 18.52 & 13.43 \\
1 lab, Qwen loop & 13.85 & 18.52 & 13.43 \\
10 labs, fixed head & 13.81 & \textbf{22.22} & \textbf{28.36} \\
10 labs, Qwen loop & 13.81 & \textbf{22.22} & \textbf{28.36} \\
\midrule
Uniform random expectation & 20.00 & 20.00 & 20.00 \\
\bottomrule
\end{tabular}
\end{table}

\begin{longtable}{lrrrrr}
\caption{Complete pilot metrics. Accuracy columns and Top-3 are percentages. MRR is higher-is-better; cross-entropy (CE) is lower-is-better.}\\
\toprule Method & Macro Top-1 & Micro Top-1 & MRR & Top-3 & CE \\ \midrule
\endfirsthead
\toprule Method & Macro Top-1 & Micro Top-1 & MRR & Top-3 & CE \\ \midrule
\endhead
\multicolumn{6}{l}{\textit{VCC: 20 interventions, 360 queries}} \\
Public head + cosine & \textbf{14.03} & \textbf{15.28} & 0.4087 & 52.50 & \textbf{1.6365} \\
1 lab, fixed head & 13.14 & 13.06 & 0.4004 & 56.39 & 1.6537 \\
1 lab, Qwen direct & 13.14 & 13.06 & 0.4004 & 56.39 & 1.6537 \\
1 lab, Qwen loop & 13.85 & 14.44 & \textbf{0.4102} & \textbf{56.94} & 1.6464 \\
10 labs, fixed head & 13.81 & 13.33 & 0.4010 & 52.50 & 1.6672 \\
10 labs, Qwen loop & 13.81 & 13.33 & 0.4006 & 52.50 & 1.6675 \\
\midrule
\multicolumn{6}{l}{\textit{Norman: 27 pairs, 1350 profiles}} \\
Public head + cosine & 17.85 & 17.85 & 0.4555 & 61.93 & 1.6204 \\
1 lab, fixed head & 18.52 & 18.52 & 0.4730 & \textbf{70.37} & 1.5824 \\
1 lab, Qwen direct & 18.52 & 18.52 & 0.4730 & \textbf{70.37} & 1.5824 \\
1 lab, Qwen loop & 18.52 & 18.52 & 0.4730 & \textbf{70.37} & 1.5824 \\
10 labs, fixed head & \textbf{22.22} & \textbf{22.22} & 0.4768 & 67.78 & \textbf{1.5822} \\
10 labs, Qwen loop & \textbf{22.22} & \textbf{22.22} & \textbf{0.4787} & 69.41 & 1.5832 \\
\midrule
\multicolumn{6}{l}{\textit{Tahoe: 67 drugs, 67 queries}} \\
Public head + cosine & 20.90 & 20.90 & 0.4669 & 64.18 & 1.6094 \\
1 lab, fixed head & 13.43 & 13.43 & 0.4157 & 56.72 & 1.7123 \\
1 lab, Qwen direct & 13.43 & 13.43 & 0.4157 & 56.72 & 1.7123 \\
1 lab, Qwen loop & 13.43 & 13.43 & 0.4157 & 56.72 & 1.7123 \\
10 labs, fixed head & \textbf{28.36} & \textbf{28.36} & \textbf{0.5177} & 65.67 & 1.5879 \\
10 labs, Qwen loop & \textbf{28.36} & \textbf{28.36} & 0.5172 & \textbf{67.16} & \textbf{1.5865} \\
\midrule
\end{longtable}


\paragraph{Interpretation.}
The first pass completes 33 physical fits, 27 local-Qwen proposals and 18 arm-level test evaluations. On double-gene perturbations, macro Top-1 rises from 18.52\% for the single-client fixed head to 22.22\% for the ten-client loop; on Tahoe it rises from 13.43\% to 28.36\%. The corresponding fixed ten-client heads attain the same Top-1, so these gains do not establish a benefit from loop organization. VCC remains below chance: 13.14\% for the single-client fixed head, 13.85\% for the single-client loop and 13.81\% for the ten-client loop. The public head reaches 14.03\%, so supervised head adaptation does not improve the primary VCC metric in this pilot.

Intervention-cluster paired bootstrap intervals for ten-client-loop minus one-client-fixed macro Top-1 are $[-2.15,3.66]$, $[-11.33,18.74]$, and $[1.49,28.36]$ percentage points for VCC, Norman and Tahoe, respectively (10,000 resamples). These exploratory intervals do not replace independent training seeds, client rotations, or multiplicity-adjusted confirmatory tests. The very small local development sets, frozen backbone and short schedule are material limitations. The pilots establish executable training and evaluation across the requested strata, not SOTA performance, full-model adaptation or a completed three-task uniform federated benchmark spanning DTI and proteomics.

\paragraph{Reproduction records.}
The local extension \texttt{extensions/federated\_bio\_20260918} contains pinned-backbone loading, three data adapters, local workers, Qwen proposals, training, selection sealing and an independent result audit. Data manifests record source identities, split memberships, candidate rosters and checkpoint hashes. All 18 sets of scalar metrics are reconstructed from saved ranks; client/split disjointness and true-option inclusion are checked. Raw biological arrays, hidden-state caches and model checkpoints are not included in the manuscript repository. Machine-readable aggregate results accompany the generated tables.
